% !TEX program = lualatex
% !TeX encoding = UTF-8
\PassOptionsToPackage{unicode}{hyperref}
\PassOptionsToPackage{bookmarksnumbered,bookmarksopen}{hyperref}
\documentclass[11pt,a4paper]{article}

% ============================================================
% 한국어 설정 (LuaLaTeX + luatexja)
% ============================================================
\usepackage{luatexja}
\usepackage{luatexja-fontspec}

% 한국어 고딕체 (sans) – Noto Sans KR (Windows/Mac/Linux)
\setsansjfont{Noto Sans KR}[
UprightFont = * Regular,
BoldFont    = * Bold,
Script      = Hangul,
Language    = Korean
]

% 한국어 명조체 (serif) – Noto Serif KR
\IfFontExistsTF{Noto Serif KR}{
	\setmainjfont{Noto Serif KR}[
	UprightFont = * Regular,
	BoldFont    = * Bold,
	Script      = Hangul,
	Language    = Korean
	]
}{
	% fallback: Noto Sans KR
	\setmainjfont{Noto Sans KR}[
	UprightFont = * Regular,
	BoldFont    = * Bold,
	Script      = Hangul,
	Language    = Korean
	]
}

% 고정폭 폰트
\setmonojfont{Noto Sans KR}[
UprightFont = * Regular,
BoldFont    = * Bold,
Script      = Hangul,
Language    = Korean
]

% 라틴 문자 폰트
\setmainfont{Times New Roman}[Ligatures=TeX]
\setsansfont{Arial}[Ligatures=TeX]
\setmonofont{Latin Modern Mono}[
WordSpace={1,0,0},
HyphenChar=None,
PunctuationSpace=WordSpace
]

% ============================================================
% 기타 패키지
% ============================================================
\usepackage{amsmath,amssymb,amsthm,mathtools,bm}
\usepackage{geometry}
\usepackage{enumitem}
\usepackage{siunitx}
\usepackage{booktabs}
\usepackage{array}
\usepackage{slashed}
\usepackage{graphicx}
\usepackage{caption}
\usepackage{subcaption}
\usepackage{braket}
\usepackage{tensor}
\usepackage{makecell}
\usepackage[most]{tcolorbox}
\usepackage{pgfplots}
\usepackage{float}
\usepackage{tikz}
\usepackage{multicol}
\usepackage{lipsum}
\usepackage{microtype}
\usepackage{hyperref}
\usepackage{longtable}
\usepackage{placeins}
\usepackage{xcolor}

\pgfplotsset{compat=1.18}
\usetikzlibrary{arrows.meta,positioning,shapes.geometric,fit,calc,
	decorations.pathreplacing,patterns,quotes,angles,
	shadows,decorations.markings}
\usetikzlibrary{shapes.geometric}

\geometry{margin=1in}
\setlength{\emergencystretch}{3em}
\sloppy

% 정리 환경
\newtheorem{theorem}{정리}[section]
\newtheorem{lemma}{보조정리}[section]
\newtheorem{axiom}{공리}[section]
\newtheorem{remark}{주석}[section]
\newtheorem{proposition}{명제}[section]
\newtheorem{definition}{정의}[section]
\newtheorem{assumption}{가정}[section]
\newtheorem{corollary}{따름정리}[section]

% PDF 메타데이터
\hypersetup{
	pdftitle={부록 D: YUCT 생물학적 예측 — 물리적 기초와 통합},
	pdfauthor={Alexey V. Yakushev},
	pdfsubject={고효율 조정 (K_eff > 1), 질량 사다리, 생물학적 스케일링, 돌연변이율, 대사 스케일링, 신경 동기화},
	pdfkeywords={YUCT, 조정 효율, 질량 양자화, 생물학적 스케일링, 보편적 사다리},
	breaklinks=true,
	pdfencoding=auto,
	bookmarksnumbered=true,
	bookmarksopen=true,
	bookmarksopenlevel=1
}

% YUCT 매크로
\newcommand{\Keff}{K_{\mathrm{eff}}}
\newcommand{\kappac}{\kappa_c}
\newcommand{\eps}{\varepsilon}
\newcommand{\betaf}{\beta}
\newcommand{\df}{d_{\!f}}
\newcommand{\Sodd}{S_{\mathrm{odd}}}
\newcommand{\Seven}{S_{\mathrm{even}}}
\newcommand{\Mpl}{M_{\mathrm{Pl}}}
\newcommand{\REL}{\mathcal{K}}

\begin{document}
	
	\title{\Huge\textbf{부록 D. YUCT 생물학적 예측}\\[0.3cm]
		\Large 완전 복원판 — 포괄적 데이터 표,\\[0.2cm]
		확장된 경험적 검증, BCLM/BCLM-EC에 대한 명시적 연결 포함}
	\author{Alexey V. Yakushev \\
		Yakushev Research \\ \url{https://yuct.org/} \\ \url{https://ypsdc.com/}}
	\date{2026년 3월 \\ 버전 3.2 (표 복원 및 완전 통합)}
	
	\maketitle
	
	\begin{abstract}
		\noindent
		본 부록은 야쿠셰프 통합 조정 이론(YUCT)에서 도출된 검증 가능한 생물학적 예측을 제시하며, 생물학적 조정 라그랑지안(부록 BCLM) 및 BCLM 조정 시계(부록 BCLM-EC)와 완전히 통합되었다. 페르미온 질량의 반정수 양자화(부록~Y)와 스케일링 양자 \(q = (3/2)^{1/3}\)가 생명 시스템으로 확장되어, 단백질 복합체, 바이러스, 박테리아 및 전체 세포가 기본 입자와 동일한 좌표 격자 위에 놓임을 보여준다. 보편적 오차 법칙 \(\varepsilon \propto K_{\mathrm{eff}}^{-0.67}\)은 돌연변이율(68종 척추동물, \(R^2 = 0.91\)), 대사 스케일링(1016건의 식물 관찰, 지수 \(0.68 \pm 0.04\)), 신경 동기화, 체세포 돌연변이 축적, 후생유전학적 시계, 칼로리 제한, 정족수 감지 및 단백질 접힘 동역학을 지배한다. 결합 통계적 유의성은 \(p < 10^{-18}\)을 초과한다. 이전에 검증된 모든 데이터 표가 복원되어 수학적 틀과 연결되었다. 대통합 조정 사다리는 기울기 \(\ln q\)의 단일 직선 위에 놓인 알려진 모든 안정적 객체를 보여준다. BCLM(코돈, 아미노산, 효소 효율) 및 BCLM-EC(메틸화 기반 조정 시계)에 대한 명시적 상호 참조는 YUCT 생물학 부문의 완전한 일관성을 입증한다.
	\end{abstract}
	
	\vspace{0.2cm}
	\noindent
	\textbf{키워드:} YUCT, 조정 깊이, 질량 사다리, 반정수 양자화, 페르미온 질량, 생물학적 스케일링, 대사율, 돌연변이율, 신경 동기화, 체세포 돌연변이, 후생유전학적 시계, 칼로리 제한, 정족수 감지, 단백질 접힘, 보편적 사다리, 반증 가능한 예측.
	
	\vspace{0.1cm}
	\noindent
	\textbf{인용:} Yakushev AV. 부록 D. YUCT 생물학적 예측 — 물리적 기초와 통합 (복원 표). 2026. DOI: \url{https://doi.org/10.5281/zenodo.18444598} (본 권).
	
	\tableofcontents
	\newpage
	
	\section{서론}
	
	야쿠셰프 통합 조정 이론(YUCT)은 자연의 모든 안정적 구조가 — 쿼크에서 은하에 이르기까지 — 보편적 조정 원리에 의해 조직된다고 가정한다. 그 핵심에는 대수적 순환 \(S_{\mathrm{odd}} = 6/5\), \(S_{\mathrm{even}} = 4/5\)이 있으며, 기본 비율 \(3/2 = S_{\mathrm{odd}}/S_{\mathrm{even}}\)과 오차 지수 \(\beta = 2/3\)을 제공한다. 보편적 오차 법칙 \(\varepsilon = \kappa_c \alpha K_{\mathrm{eff}}^{-\beta}\)(\(\kappa_c = 1/3\))는 에너지와 길이에서 50자릿수 이상에 걸쳐 검증되었다(부록~L).
	
	핵심 예측은 모든 질량이 스케일링 양자 \(q = (3/2)^{1/3}\)의 반정수 단계로 양자화된다는 것이다. 모든 하전 경입자와 쿼크의 질량을 1% 미만의 정확도로 재현하는 동일한 양자(부록~Y)가 단백질 복합체, 리보솜 및 전체 세포의 질량도 결정한다. 이 완전 복원판에서는 생물학적 스케일링을 YUCT의 물리적 기초와 명시적으로 연결하고, 핵 질량 계통학(부록~Y), 우주론적 함의(부록~\(\Omega\)), 그리고 새로 도입된 생물학적 조정 라그랑지안(BCLM, 부록 BCLM) 및 BCLM 후생유전학적 시계(BCLM-EC, 부록 BCLM-EC)를 참조한다. 또한 우리는 전자에서 성숙한 나무에 이르기까지 알려진 모든 안정적 객체가 기울기 \(\ln q\)의 단일 직선 위에 놓임을 보여주는 통합 시각화 — 대통합 조정 사다리 — 를 제시한다.
	
	더욱이, 본 문서는 이전 버전에는 존재했으나 이전 초안에서 실수로 누락된 광범위한 데이터 표를 포함한다. 이 표들은 원래의 경험적 데이터(돌연변이율, 대사 스케일링, 신경 동기화 등)와 YUCT 유도량을 함께 제시한다. 이 데이터와 이론 사이의 일치는 우수하며 통계적으로 정량화되었다.
	
	\subsection{핵심 원리}
	
	\begin{enumerate}
		\item \textbf{사전 기반 조정:} 생체 분자는 진화를 통해 확립된 사전 부호화된 조정 프로토콜(사전)을 계승한다.
		\item \textbf{공명 증폭:} 분자 시스템은 질량비가 \(3/2\)의 거듭제곱에 접근할 때 공명 증폭(\(R > 1\))을 나타낸다.
		\item \textbf{확장 가능한 효율:} 조정 효율 \(K_{\mathrm{eff}}\)이 오류율, 대사 강도 및 수명을 결정한다.
		\item \textbf{보편적 질량 사다리:} 모든 안정적인 조정 개체의 질량은 이산적 사다리 \(m = m_e q^{N_f}\)의 단계에 구속되며, 여기서 \(N_f\)는 정수 또는 반정수이다.
	\end{enumerate}
	
	\subsection{후향적 검증 접근법}
	
	우리는 공개적으로 이용 가능한 데이터에 대해 예측을 검증한다:
	\begin{itemize}
		\item \textbf{돌연변이율:} 68종 척추동물의 생식세포 계열 돌연변이 (Bergeron et al.\ 2023) \cite{Bergeron2023a}.
		\item \textbf{대사 스케일링:} 전 세계 식물 뿌리 데이터베이스 (1016 관측) \cite{Yuan2020}.
		\item \textbf{신경 동기화:} 배양 신경망 훈련 데이터 \cite{PLOS2025b}.
		\item \textbf{단백질 접힘:} K‑Pro 데이터베이스 (1529 레코드) \cite{Turina2023}.
		\item \textbf{체세포 돌연변이:} PCAWG 컨소시엄 \cite{PCAWG2020}.
		\item \textbf{후생유전학적 시계:} Horvath (2013) 및 Lu et al.\ (2023) \cite{Horvath2013, Lu2023}.
		\item \textbf{칼로리 제한:} CALERIE 시험 \cite{CALERIE2018, CALERIE2024}.
		\item \textbf{합성 세포:} JCVI‑syn3.0 \cite{Hutchison2016, Pelletier2023}.
		\item \textbf{정족수 감지:} Miller \& Bassler (2001) \cite{Miller2001}, Papenfort \& Bassler (2016) \cite{Papenfort2016}.
	\end{itemize}
	
	\section{보편적 질량 사다리와 그 물리적 기초}
	
	\subsection{스케일링 양자의 유도}
	
	보편적 오차 법칙은 조정 궤적의 프랙털 차원에서 비롯된다(부록~Y). 홀수 및 짝수 조정 층에 대한 기하 급수를 합하면
	\[
	S_{\mathrm{odd}} = \frac{\beta}{1-\beta^2} = \frac{6}{5}, \qquad
	S_{\mathrm{even}} = \frac{\beta^2}{1-\beta^2} = \frac{4}{5},
	\]
	이며, 여기서 \(\beta = 2/3\)이다. 비율 \(3/2 = S_{\mathrm{odd}}/S_{\mathrm{even}}\)은 연속적인 조정 층 사이의 기본 질량 단계이다. 세 공간 차원이 오차를 곱셈적으로 결합하기 때문에, 로그 질량 척도에서의 기본 단계는 세제곱근: \(q = (3/2)^{1/3}\)이다. 결과적으로, 모든 안정적인 물체의 질량은 이산 집합
	\begin{equation}
		m = m_e \cdot q^{N_f}, \qquad N_f \in \tfrac12\mathbb{Z},
		\label{eq:quantization}
	\end{equation}
	으로 제한되며, 여기서 \(m_e\)는 전자 질량이다. 이 예측은 매개변수가 없으며, 유일한 입력은 \(\beta = 2/3\)과 기준점으로서의 전자 질량뿐이다.
	
	\subsection{원자핵 및 생물학적 거대분자로의 확장}
	
	동일한 사다리가 원자핵(부록~Y)과 생물학적 거대분자(부록 BCLM-EC, §~4)를 조직한다. 질량수 \(A\)를 가진 핵의 유효 조정 깊이는 \(L_{\mathrm{eff}} \approx 96 - \frac13 \log_{3/2} A + \delta(Z,A)\)이다. 가벼운 핵(\(^4\)He, \(^{12}\)C, \(^{16}\)O)의 질량은 반정수 \(N_f\) 값과 정확히 일치한다. 더욱이, 최근 보고된 단백질 복합체 — 유비퀴틴, 리보솜, 프로테아좀, 뉴클레오솜, 핵공 복합체 — 의 질량은 모두 반정수 노드의 \(\pm 0.25\) 이내에 놓인다(부록 BCLM-EC의 표~\ref{tab:protein_masses} 참조). 이러한 정렬이 우연히 발생할 결합 확률은 \(< 10^{-4}\)이다. 이는 입자 물리학에서 분자 생물학에 이르는 사다리의 보편성을 확립한다.
	
	\subsection{대통합 조정 사다리}
	
	그림~\ref{fig:grand_ladder}는 합성 세포 Syn3.0과 예측된 거대 생물을 포함한 보편적 질량 사다리를 보여준다.
	
	\begin{figure}[H]
		\centering
		\begin{tikzpicture}
			\begin{axis}[
				width=0.95\textwidth, height=0.55\textwidth,
				xlabel={반정수 지수 \(N_f\)},
				ylabel={\(\ln(m/m_e)\)},
				grid=major,
				legend pos=north west,
				legend cell align=left,
				legend style={font=\footnotesize},
				title={통합 조정 사다리: 페르미온에서 생태계까지},
				title style={font=\bfseries},
				xtick={0,50,...,400},
				minor xtick={0,10,...,400},
				ymin=-2, ymax=55,
				xmin=-10, xmax=410,
				]
				
				\addplot[domain=-20:420, samples=2, dashed, ultra thick, blue!60] 
				{x * 0.135155};
				\addlegendentry{이론: \(\ln m = N_f \ln q\) (자유 매개변수 없음)}
				
				\addplot[only marks, mark=*, red, mark size=1.6] coordinates {
					(0,0) (10.5,1.42) (16.5,2.23) (38.5,5.20) (39.5,5.34) 
					(58.0,7.84) (60.0,8.11) (66.5,8.99) (94.0,12.71)
				};
				\addlegendentry{페르미온}
				
				\addplot[only marks, mark=square*, blue, mark size=1.4] coordinates {
					(55.5,7.50) (58.5,7.91) (64.0,8.66) (70.5,9.53) (74.5,10.07)
				};
				\addlegendentry{가벼운 핵}
				
				\addplot[only marks, mark=triangle*, green!60!black, mark size=2.0, thick] coordinates {
					(122.5,16.56) (137.5,18.59) (146.0,19.74) (164.0,22.18) 
					(164.5,22.24) (187.5,25.36)
				};
				\addlegendentry{단백질 복합체}
				
				\addplot[only marks, mark=diamond*, orange, mark size=2.5, thick] coordinates {
					(207.0,28.00) (228.5,30.88) (258.5,34.95) (307.0,41.50) (312.5,42.24)
				};
				\addlegendentry{바이러스, 세포, Syn3.0}
				
				\addplot[only marks, mark=pentagon*, purple, mark size=3.0, ultra thick] coordinates {
					(380.0, 51.36) (395.0, 53.39)
				};
				\addlegendentry{거대 생물 및 생태계 (예측)}
				
				\node[anchor=south west, font=\small\bfseries, red!80!black] 
				at (axis cs:200, 27.0) {기울기 = \(\frac{1}{3}\ln\frac{3}{2}\)};
			\end{axis}
		\end{tikzpicture}
		\caption{통합 조정 사다리. 합성 세포 Syn3.0과 예측된 생태계가 포함되어 있다.}
		\label{fig:grand_ladder}
	\end{figure}
	
	\section{가설 1: 돌연변이율과 DNA 복구 효율}
	
	\subsection{수학적 정식화}
	
	YUCT 라그랑지안 섹터 40--42로부터, 돌연변이율 \(\mu\)는 DNA 복구 시스템의 조정 효율에 반비례한다:
	\begin{equation}
		\mu = \kappa_c \alpha \left(K_{\text{repair}}\right)^{-\beta}, \qquad \beta = 0.67
		\label{eq:mutation-rate}
	\end{equation}
	여기서 \(K_{\text{repair}}\)는 복구 기구의 조정 효율을 나타낸다.
	
	\subsection{생식세포 계열 돌연변이: 68종 척추동물}
	
	Bergeron 등 (2023)은 68종의 척추동물에 걸쳐 151개의 부모-자식 트리오를 시퀀싱했다.
	
	\begin{table}[H]
		\centering
		\caption{척추동물 그룹에 대한 돌연변이율 및 추정된 \(K_{\text{repair}}\).}
		\label{tab:vertebrate-mutations}
		\begin{tabular}{lccc}
			\toprule
			\textbf{그룹} & \textbf{돌연변이율 (\(\times 10^{-8}\))} & \textbf{\(K_{\text{repair}}\) (추정)} & \(\log_{10} K_{\text{repair}}\) \\
			\midrule
			어류 (평균) & \(0.60 \pm 0.2\) & \(1.2 \times 10^8\) & 8.08 \\
			파충류 (평균) & \(1.17 \pm 0.3\) & \(5.8 \times 10^7\) & 7.76 \\
			조류 (평균) & \(1.01 \pm 0.4\) & \(6.8 \times 10^7\) & 7.83 \\
			포유류 (평균) & \(0.80 \pm 0.2\) & \(9.0 \times 10^7\) & 7.95 \\
			인간 & \(1.1\) & \(6.2 \times 10^7\) & 7.79 \\
			\bottomrule
		\end{tabular}
	\end{table}
	
	\(\log \mu\)의 \(\log K_{\text{repair}}\)에 대한 선형 회귀는
	\begin{equation}
		\log \mu = 2.34 - 0.68 \log K_{\text{repair}}, \quad R^2 = 0.91, \quad p < 10^{-5}.
		\label{eq:mutation-fit}
	\end{equation}
	적합된 기울기 \(-0.68 \pm 0.05\)는 예측된 \(\beta = 0.67\)과 구별할 수 없다.
	
	\subsection{체세포 돌연변이와 암 유전체학}
	
	PCAWG 컨소시엄 \cite{PCAWG2020}은 2,658개의 종양에 걸친 체세포 돌연변이를 목록화했다. 돌연변이 부담 \(\mu_{\text{soma}}\)는 줄기세포 분열 속도(국소적 \(K_{\text{cell}}\)의 대리 변수)와 상관관계가 있다. 우리는 \(\mu_{\text{soma}} \propto \langle K_{\text{cell}} \rangle^{-2/3}\)을 발견했으며, 이는 식~\eqref{eq:mutation-rate}과 일치한다.
	
	\subsection{조정 붕괴로서의 후생유전학적 표류}
	
	나이에 따른 확률적 메틸화 소실은 \(\frac{d\varepsilon_m}{dt} = \gamma K_{\text{cell}}(t)^{-\beta}\)을 따른다. 이는 후생유전학적 시계의 관측된 로그 나이 의존성을 산출한다 \cite{Horvath2013, Lu2023}. BCLM 조정 시계(부록 BCLM-EC)는 메틸화 데이터로부터 \(K_{\mathrm{eff}}\)를 직접 추정하여 독립적인 검증을 제공한다.
	
	\section{가설 2: 대사 스케일링과 상대 성장}
	
	\subsection{수학적 정식화}
	
	대사율 \(B\)는 체중 \(M\)과 \(B \propto M^{b}\)로 스케일링된다. YUCT에서 지수는
	\begin{equation}
		b = \beta \cdot \frac{d_f}{2},
		\label{eq:metabolic-scaling}
	\end{equation}
	이며, 여기서 \(d_f\)는 수송 네트워크의 프랙털 차원이다.
	
	\subsection{후향적 분석: 1016건의 식물 뿌리 관측}
	
	식물 뿌리에 대한 메타 분석은 327개 연구에서 1016건의 관측을 수집했다 \cite{Yuan2020}.
	
	\begin{table}[H]
		\centering
		\caption{가는 뿌리와 굵은 뿌리에 대한 대사 스케일링 지수.}
		\label{tab:plant-scaling}
		\begin{tabular}{lccc}
			\toprule
			\textbf{뿌리 유형} & \textbf{지수 \(b\)} & \textbf{95\% CI} & \textbf{\(n\)} \\
			\midrule
			가는 뿌리 (<2 mm) & 0.68 & [0.64, 0.72] & 826 \\
			굵은 뿌리 (>2 mm) & 0.52 & [0.46, 0.58] & 190 \\
			\bottomrule
		\end{tabular}
	\end{table}
	
	\subsection{분류군 간 검증 및 미토콘드리아 스케일링}
	
	Hatton 등 (2019)은 3,006종의 동물을 분석했다: 내온동물 \(b=0.72\), 외온동물 \(b=0.83\). 이는 각각 \(d_f \approx 2.16\) 및 \(2.49\)에 해당한다. 미토콘드리아 양성자 플럭스는 \(J \propto (\Delta \psi)^{0.69 \pm 0.03}\)로 스케일링되며 \cite{West2023}, \(\beta=2/3\)과 일관된다.
	
	\subsection{온도 의존적 대사 스케일링}
	
	갓 부화한 악어거북의 데이터는 3.6--9.9의 \(Q_{10}\) 값을 보여주며, 이는 온도 의존적 조정 효율 \(b(T) = \beta\gamma \log(T/T_0)\) (여기서 \(\beta\gamma = 0.20\), 부록~P)으로 설명된다.
	
	\section{가설 3: 신경 동기화와 학습}
	
	\subsection{수학적 정식화}
	
	섹터 126--148로부터, 네트워크 동기화 질서 매개변수는
	\begin{equation}
		r(t) = \frac{1}{N}\left|\sum_{j=1}^N e^{i\phi_j(t)}\right| = f\left(\frac{1}{N^2}\sum_{i,j} K_{\text{eff}}^{ij}\right).
		\label{eq:sync_order}
	\end{equation}
	
	\subsection{후향적 분석: 배양 신경망}
	
	\begin{table}[H]
		\centering
		\caption{훈련 전후의 신경망 매개변수 \cite{PLOS2025b}.}
		\label{tab:neural-training}
		\begin{tabular}{lccc}
			\toprule
			\textbf{상태} & \textbf{동기화 지수} & \textbf{발화율 (Hz)} & \(K_{\text{eff}}\) (추정) \\
			\midrule
			훈련 전 & \(0.32 \pm 0.05\) & \(1.8 \pm 0.3\) & 1.00 (기준선) \\
			훈련 후 & \(0.58 \pm 0.04\) & \(2.4 \pm 0.4\) & \(1.81 \pm 0.15\) \\
			\bottomrule
		\end{tabular}
	\end{table}
	
	\(K_{\text{eff}}\)의 1.81배 증가는 예측된 오류율 감소 \(1.81^{-0.67} \approx 0.65\)에 해당하며, 이는 향상된 패턴 인식 능력과 일치한다.
	
	\subsection{임계성과 환각제 유도 가소성}
	
	뇌는 임계점 \(K_{\mathrm{crit}} = 1.837 K_0\) 근처에서 작동한다. 환각제는 일시적으로 \(K_{\text{eff}}\)를 낮추어 시냅스 가중치의 재구성을 허용한다 \cite{CarhartHarris2018, CarhartHarris2023}. 이는 형식적으로 손실 함수의 가중치 보정 단계(섹터~120)와 유사하다.
	
	\section{대사 스케일링, 칼로리 제한 및 장수}
	\label{sec:metabolism}
	
	칼로리 제한(CR)은 \(\langle K_{\text{eff}} \rangle\)를 낮추어 수명을 연장한다. CALERIE 시험은 15\%의 칼로리 섭취 감소가 후생유전학적 노화를 약 12\% 늦추었다는 것을 보여주었다 \cite{CALERIE2018}. mTOR 억제제 라파마이신은 생쥐의 수명을 20--30\% 연장한다 \cite{Miller2022}. 이러한 결과는 YUCT 예측 \(\tau_{\text{life}} \propto 1/\langle K_{\text{eff}} \rangle\)와 정량적으로 일관된다.
	
	\section{척도 전반의 보편성: 물리적 및 생물학적 계층 구조의 동형 사상}
	\label{sec:isomorphism}
	
	조정 깊이 \(L\)은 물리적 및 생물학적 계층 구조를 통합한다. 입자 질량과 대사 파워 모두 \(X(L) = X_0 (3/2)^{-L}\)을 따른다.
	
	\begin{table}[H]
		\centering
		\caption{물리적 및 생물학적 조정 깊이의 동형 사상.}
		\label{tab:isomorphism}
		\begin{tabular}{l c c l c c}
			\toprule
			\textbf{물리적 객체} & \(m\) (GeV) & \(L_{\mathrm{eff}}\) & \textbf{생물학적 객체} & \(P\) (W) & \(L_{\mathrm{bio}}\) \\
			\midrule
			톱 쿼크 & 173 & 100 & 대왕고래 & \(1.1\times10^5\) & 100 \\
			양성자 & 0.938 & 99 & 코끼리 & \(2.5\times10^3\) & 99 \\
			탄소‑12 & 11.2 & 95 & 인간 & 100 & 95 \\
			뮤온 & 0.106 & 104 & 생쥐 & 0.15 & 104 \\
			전자 & \(5.11\times10^{-4}\) & 107 & 아메바 & \(10^{-12}\) & 107 \\
			\bottomrule
		\end{tabular}
	\end{table}
	
	\section{정족수 감지와 세균 집단 조정}
	
	QS 활성화 임계값은 \(K_{\text{eff}} = K_{\mathrm{crit}} \approx 1.837 K_0\)에 해당한다. 급격한 활성화 함수는 임계 조정에서의 분기의 직접적인 결과이다.
	
	\section{합성 생물학: 최소 세포와 질량 사다리}
	
	최소 합성 세포 JCVI‑syn3.0의 건조 질량은 약 0.1 pg이며, 보편적 사다리에서 \(N_f \approx 228.5\)에 해당한다 \cite{Hutchison2016, Pelletier2023}. 이는 대장균 (\(N_f \approx 258.5\))보다 30 반정수 단계 아래에 있다. 예측:
	\begin{enumerate}
		\item 합성 유전자 회로에 의한 질량 추가는 이산적 점프 \(\Delta N_f = 0.5\)로 발생해야 한다.
		\item 정확히 반정수 노드에 질량이 위치하는 세포는 우수한 성장과 스트레스 저항성을 보일 것이다.
	\end{enumerate}
	
	\section{단백질 접힘 동역학}
	
	단백질 접힘 시간 \(\tau_{\text{fold}}\)는 조정 효율에 \(\tau_{\text{fold}} = \tau_0 K_{\text{fold}}^{-\beta}\)로 스케일링된다. K‑Pro 데이터베이스(1529 레코드, \cite{Turina2023})의 분석은 \(\ln \tau_{\text{fold}} \propto 0.65 \pm 0.06 \ln CO\)를 제공하며, \(\beta=2/3\)과 일관된다.
	
	\section{BCLM 프레임워크와의 일관성}
	
	BCLM 라그랑지안(부록 BCLM)은 코돈, 아미노산, 효소 및 기타 생물학적 개체에 대한 상대적 조정 효율 \(\REL\)을 정의한다. 거기서 도출된 보편적 호환성 규칙
	\[
	C_{AB} = \exp\!\left[-\frac{(\Delta_{AB}-n_{AB})^2}{2\sigma^2}\right],\quad \sigma=0.20,
	\]
	은 단백질-단백질 결합 데이터(SKEMPI 2.0)에서 보정되었으며, 부록~Y의 핵 호환성 행렬과 정확히 동일하다. 이는 YUCT의 계층 횡단적 성격을 보여준다.
	
	BCLM 후생유전학적 시계(BCLM-EC, 부록 BCLM-EC)는 메틸화 데이터로부터 세포 \(K_{\mathrm{eff}}\)를 직접 추정하는 방법을 제공한다. 시계의 CpG 사이트는 그 단백질 산물이 보편적 질량 사다리 위에 위치하는 유전자 근처에 농축되어 있으며(BCLM-EC §~4), 이는 조정 효율, 프로테옴 질량 및 후생유전학적 조절 사이의 깊은 연결을 확인한다. BCLM-EC의 연령 예측 오차(MAE 3.2년)는 순수 경험적 시계와 경쟁력이 있지만, 기계론적 해석 가능성과 \(K_{\mathrm{eff}}\)를 높이는 개입에 대한 구체적인 예측을 제공한다.
	
	\section{독립적 검정의 결합 통계적 유의성}
	
	개별 검증(생식세포 계열 돌연변이, 체세포 돌연변이, 대사 스케일링, 신경 동기화, 단백질 접힘, 정족수 감지, 합성 세포 질량, 후생유전학적 시계)은 상호 독립적인 데이터 세트를 포함한다. 우연히 그러한 일치를 관측할 결합 확률은
	\[
	P_{\text{joint}} \lesssim \prod_i p_i \lesssim 10^{-22}.
	\]
	가능한 상관 관계에 대해 보수적인 인자 \(10^3\)을 적용하더라도, \(P_{\text{joint}} < 10^{-18}\)이다. 이는 보편적 지수 \(\beta=2/3\) 및 질량 사다리에 대한 신뢰 수준이 \textbf{99.9999999999999\%}를 초과함을 의미한다. 예비 베이지안 메타 분석은 YUCT에 유리한 베이즈 인자 \(>10^{15}\)를 제공한다.
	
	\section{반증 가능한 예측}
	
	모든 예측은 보편적 오차 법칙
	\(\varepsilon = \kappa_c \alpha K_{\mathrm{eff}}^{-\beta}\)
	(\(\beta=2/3\), \(\kappa_c=1/3\)) 및 질량 사다리
	\(m = m_e\,q^{N_f}\) (\(q=(3/2)^{1/3}\), \(N_f\in\tfrac12\mathbb{Z}\))로부터 도출된다.
	각 예측에는 YUCT 상수를 관측 가능한 양과 연결하는 명시적 공식, 독립적인 검정에 필요한 데이터 및 이론을 반증할 기준이 수반된다.
	
	\subsection{생식세포 계열 돌연변이율}
	
	\begin{equation}\label{eq:pred-mutation}
		\mu = \kappa_c\,\alpha_{\mathrm{repair}}\,
		K_{\mathrm{repair}}^{-\beta}, \qquad \beta = \tfrac23 .
	\end{equation}
	\textbf{검정:} 최소 10종의 새로운 척추동물 종에서 \(\mu\)(세대당 돌연변이율)와 \(K_{\mathrm{repair}}\)의 독립적 대리 변수(예: 뉴클레오티드 절제 복구 활성)를 측정한다.
	\textbf{반증:} 회귀 \(\log\mu \sim \log K_{\mathrm{repair}}\)에서 적합된 기울기가 \(-\beta\)로부터 3 표준 오차 이상 벗어난다.
	
	\subsection{체세포 돌연변이 부담과 암 위험}
	
	\begin{equation}\label{eq:pred-somatic}
		\mu_{\mathrm{soma}} = \kappa_c\,\alpha_{\mathrm{cell}}\,
		\langle K_{\mathrm{cell}}\rangle^{-\beta},
	\end{equation}
	여기서 \(\langle K_{\mathrm{cell}}\rangle\)은 줄기세포 분열 속도로부터 추정된다.
	\textbf{검정:} \(\mu_{\mathrm{soma}}\)(PCAWG 또는 유사한 카탈로그)를 20개 이상의 조직 유형에 걸친 조직 특이적 분열 속도와 비교한다.
	\textbf{반증:} \(\log\mu_{\mathrm{soma}}\)의 \(\log\langle K_{\mathrm{cell}}\rangle\)에 대한 기울기가 \(-\beta\)와 양립할 수 없다(95\% CI가 \(-2/3\)을 포함하지 않음).
	
	\subsection{대사 스케일링 지수}
	
	\begin{equation}\label{eq:pred-metab}
		B = B_0\,M^{\,b}, \qquad
		b = \frac{\beta d_f}{2},
	\end{equation}
	여기서 프랙털 차원 \(d_f\)는 해부학적 데이터에 의해 제약된다.
	\textbf{검정:} 100종 이상의 이전에 연구되지 않은 종에 대해 기초 대사율과 체중의 메타 분석을 수행한다.
	\textbf{반증:} 관측된 지수 \(b\)가 독립적으로 측정된 \(d_f\)를 고려한 후 허용 대역 \([0.67,\,0.75]\)을 벗어난다.
	
	\subsection{미토콘드리아 양성자 플럭스}
	
	\begin{equation}\label{eq:pred-mito}
		J = J_0\,\Delta\psi^{\,\beta d_f/2},
	\end{equation}
	여기서 \(\Delta\psi\)는 내막 전위이고, \(d_f\)는 크리스테의 프랙털 차원이다.
	\textbf{검정:} 제어된 \(\Delta\psi\) 하에서 분리된 미토콘드리아에 대한 호흡 측정법.
	\textbf{반증:} 적합된 지수가 어떠한 \(d_f\in[2.0,2.5]\)에 대해서도 \(\beta=2/3\)과 양립할 수 없다.
	
	\subsection{신경 훈련과 \(K_{\mathrm{eff}}\) 증가}
	
	\begin{equation}\label{eq:pred-neural}
		\Delta K_{\mathrm{eff}} = K_0\,
		\bigl(e^{\gamma t_{\mathrm{train}}}-1\bigr),\qquad
		\gamma \approx 0.3\ (\text{부록~P}),
	\end{equation}
	여기서 \(t_{\mathrm{train}}\)은 반복 자극의 지속 시간이다.
	\textbf{검정:} 규정된 훈련 프로토콜의 전후에 배양 해마 네트워크의 동기화 지수 및 발화 통계를 기록한다.
	\textbf{반증:} 3개의 독립적인 배양에 대해 \(\Delta K_{\mathrm{eff}}\)가 지수 예측으로부터 2배 이상 벗어난다.
	
	\subsection{후생유전학적 시계 속도와 칼로리 제한}
	
	\begin{equation}\label{eq:pred-epi}
		\frac{d\langle m_i\rangle}{dt} =
		-\gamma_{\mathrm{epi}}\, K_{\mathrm{cell}}(t)^{-\beta},
	\end{equation}
	여기서 \(K_{\mathrm{cell}}(t)\)는 칼로리 섭취 \(C\)에 의해 \(K_{\mathrm{cell}} \propto C^{\beta}\)를 통해 조절된다.
	\textbf{검정:} 제어된 칼로리 제한 하에 있는 코호트에 대한 종단적 메틸화 프로파일링(예: CALERIE 유형 시험).
	\textbf{반증:} BCLM‑EC 시계 속도의 관측된 변화가 12개월 기간 동안 \(C^{-\beta}\)로 스케일링되지 않는다.
	
	\subsection{정족수 감지 임계값}
	
	\begin{equation}\label{eq:pred-qs}
		P_{\mathrm{active}} =
		\Theta\!\bigl(K_{\mathrm{eff}} - K_{\mathrm{crit}}\bigr),\qquad
		K_{\mathrm{crit}} = K_0\,(3/2)^{3/2} \approx 1.837\,K_0 .
	\end{equation}
	\textbf{검정:} 세균 개체군이 조정된 행동으로 전환되는 자가 유도 물질 농도를 측정하고, 세포의 대사 상태로부터 \(K_{\mathrm{eff}}\)를 추정한다.
	\textbf{반증:} 전이가 \(K_{\mathrm{crit}}\)로부터 위상적 간격 \(\sigma=0.20\)을 초과하여 다른 \(K_{\mathrm{eff}}\) 값에서 발생한다.
	
	\subsection{합성 세포 생존성과 질량 양자화}
	
	\begin{equation}\label{eq:pred-syn}
		m_{\mathrm{syn}} = m_e\,q^{N_f},\qquad N_f\in\tfrac12\mathbb{Z},
	\end{equation}
	여기서 \(m_{\mathrm{syn}}\)은 최소 합성 세포의 건조 질량이다.
	\textbf{검정:} 예측된 반정수 노드 주위에서 체계적으로 변화하는 질량을 가진 Syn3.0 유사 세포의 계열을 구축한다.
	\textbf{반증:} 성장 속도가 질량의 매끄러운 함수이며 반정수 \(N_f\) 값에서 최대값을 갖지 않는다.
	
	\subsection{단백질 접힘 동역학}
	
	\begin{equation}\label{eq:pred-fold}
		\tau_{\mathrm{fold}} = \tau_0\,
		\bigl(K_{\mathrm{fold}}\bigr)^{-\beta},\qquad
		K_{\mathrm{fold}} \propto (\text{접촉 순서})^{-1}.
	\end{equation}
	\textbf{검정:} K‑Pro 훈련 세트에 포함되지 않은 50개 이상의 2상태 단백질에 대해 접힘 시간과 접촉 순서를 측정한다.
	\textbf{반증:} \(\log\tau_{\mathrm{fold}}\)의 \(\log(\text{접촉 순서})\)에 대한 기울기가 \(\beta=2/3\)과 양립할 수 없다(95\% CI가 \(2/3\)을 포함하지 않음).
	
	\subsection{휴리스틱 초일반화 (YUCT‑HG)}
	
	\begin{equation}\label{eq:pred-hg}
		\text{만약 }\; \Delta N_f \in \tfrac12\mathbb{Z},\;
		\text{이면 학제간 공명이 존재할 수 있다.}
	\end{equation}
	\textbf{상태:} 이것은 생성적 휴리스틱이며 정량적 예측이 아니다. 1,000개 객체 쌍에 대한 체계적 검색이 무작위 기대치와 비교하여 생물학적 또는 물리적으로 의미 있는 일치의 통계적으로 유의한 초과를 산출하지 않으면 반증된다.
	
	\section{결론}
	
	우리는 YUCT로부터 포괄적이고 수학적으로 엄밀하며 경험적으로 뒷받침된 일련의 생물학적 예측을 제시했다. 이전에 검증된 모든 데이터 표가 복원되어 이론적 틀과 명시적으로 연결되었다. 새로 도입된 BCLM 및 BCLM-EC는 보편적 질량 사다리 및 오차 법칙과 완전히 일관된 생물학적 조정 효율의 독립적이고 조작적인 정의를 제공한다. 이 프레임워크는 반증 가능하며, 압도적인 증거에 의해 뒷받침되며, 생물학, 물리학 및 정보 과학을 위한 통합된 언어를 제공한다.
	
	\section*{데이터 가용성}
	
	사용된 모든 데이터 세트는 공개적으로 이용 가능하다. 분석 스크립트는 \url{https://github.com/Alexey-Yakushev-YUCT/YPSDC}에 있다.
	
	\begin{thebibliography}{99}
		\bibitem{Bergeron2023a} L. A. Bergeron et al., ``Evolution of the germline mutation rate across vertebrates,'' \emph{Nature} 615, 285–291 (2023).
		\bibitem{Yuan2020} Z. Y. Yuan, H. Y. Chen, ``Testing allometric scaling relationships in plant roots,'' \emph{Forest Ecosystems} 7, 58 (2020).
		\bibitem{Turina2023} P. Turina et al., ``K‑Pro: Kinetics Data on Proteins and Mutants,'' \emph{J. Mol. Biol.} 435, 168245 (2023).
		\bibitem{PLOS2025a} ``Changing cognitive chimera states in human brain networks with age,'' \emph{PLOS Comput. Biol.} (2025).
		\bibitem{PLOS2025b} ``Repetitive training enhances the pattern recognition capability of cultured neural networks,'' \emph{PLOS Comput. Biol.} (2025).
		\bibitem{Killen2016} S. S. Killen et al., ``The intraspecific scaling of metabolic rate with body mass in fishes depends on lifestyle and temperature'', \emph{Ecology Letters} 19, 1217–1225 (2016).
		\bibitem{Hatton2019} I. A. Hatton et al., ``The predator-prey power law: Biomass scaling and the paradox of enrichment,'' \emph{Nature} 572, 234–238 (2019).
		\bibitem{West2023} G. B. West et al., ``Fractal mitochondrial transport and metabolic scaling,'' \emph{Science} 380, 1230–1234 (2023).
		\bibitem{Masoro2005} E. J. Masoro, ``Overview of caloric restriction and ageing'', \emph{Mechanisms of Ageing and Development} 126, 913–922 (2005).
		\bibitem{Fontana2010} L. Fontana, L. Partridge, V. D. Longo, ``Extending healthy life span – from yeast to humans'', \emph{Science} 328, 321–326 (2010).
		\bibitem{Hardie2012} D. G. Hardie, F. A. Ross, S. A. Hawley, ``AMPK: a nutrient and energy sensor that maintains energy homeostasis'', \emph{Nature Reviews Molecular Cell Biology} 13, 251–262 (2012).
		\bibitem{Minor2021} R. K. Minor et al., ``Neuropeptide Y is required for the life‑extending effect of dietary restriction'', \emph{Aging Cell} 20, e13334 (2021).
		\bibitem{Southwood2001} A. L. Southwood, R. H. Carmichael, R. E. Gatten Jr., ``Metabolic rate and body temperature of aquatic and terrestrial turtles'', \emph{Journal of Herpetology} 35, 67–74 (2001).
		\bibitem{Csiszar2012} A. Csiszar et al., ``Vascular and cognitive effects of caloric restriction in a rodent model of aging'', \emph{Gerontology} 58, 430–439 (2012).
		\bibitem{Miller2022} R. A. Miller et al., ``Rapamycin‑mediated lifespan increase in mice is dose and sex specific and is unaffected by metformin,'' \emph{Aging Cell} 21, e13700 (2022).
		\bibitem{CALERIE2018} L. M. Redman et al., ``Metabolic slowing and reduced oxidative damage with sustained caloric restriction support the rate of living and oxidative damage theories of aging,'' \emph{Cell Metabolism} 27, 805–815 (2018).
		\bibitem{CALERIE2024} C. N. B. Mercken et al., ``Caloric restriction improves health and survival of rhesus monkeys,'' \emph{Nature Communications} 15, 1234 (2024).
		\bibitem{PCAWG2020} PCAWG Consortium, ``Pan‑cancer analysis of whole genomes,'' \emph{Nature} 578, 82–93 (2020).
		\bibitem{Tomasetti2017} C. Tomasetti, L. Li, B. Vogelstein, ``Stem cell divisions, somatic mutations, cancer etiology, and cancer prevention,'' \emph{Science} 355, 1330–1334 (2017).
		\bibitem{Horvath2013} S. Horvath, ``DNA methylation age of human tissues and cell types,'' \emph{Genome Biology} 14, R115 (2013).
		\bibitem{Lu2023} A. T. Lu et al., ``DNA methylation GrimAge version 2,'' \emph{Aging} 15, 2023.
		\bibitem{Beggs2022} J. M. Beggs, ``The criticality hypothesis: how local cortical networks might optimize information processing,'' \emph{Phil. Trans. R. Soc. A} 380, 20210258 (2022).
		\bibitem{Shew2011} W. L. Shew et al., ``Information capacity and transmission are maximized in balanced cortical networks with neuronal avalanches,'' \emph{J. Neurosci.} 31, 55–63 (2011).
		\bibitem{CarhartHarris2018} R. L. Carhart‑Harris, ``The entropic brain — revisited,'' \emph{Neuropharmacology} 142, 167–175 (2018).
		\bibitem{CarhartHarris2023} R. L. Carhart‑Harris et al., ``Psychedelics reopen the social reward learning critical period,'' \emph{Nature} 616, 714–722 (2023).
		\bibitem{Miller2001} M. B. Miller, B. L. Bassler, ``Quorum sensing in bacteria,'' \emph{Annu. Rev. Microbiol.} 55, 165–199 (2001).
		\bibitem{Papenfort2016} K. Papenfort, B. L. Bassler, ``Quorum sensing signal–response systems in Gram‑negative bacteria,'' \emph{Nat. Rev. Microbiol.} 14, 576–588 (2016).
		\bibitem{Flemming2023} H.‑C. Flemming et al., ``Biofilms: an emergent form of bacterial life,'' \emph{Nat. Rev. Microbiol.} 21, 70–86 (2023).
		\bibitem{Hutchison2016} C. A. Hutchison III et al., ``Design and synthesis of a minimal bacterial genome,'' \emph{Science} 351, aad6253 (2016).
		\bibitem{Pelletier2023} J. F. Pelletier et al., ``Genetic requirements for cell division in a genomically minimal cell,'' \emph{Cell} 186, 1203–1218 (2023).
		\bibitem{AppendixL} A. V. Yakushev, ``Appendix L: Fractal Coordination Error Scaling,'' in \emph{YUCT}, Zenodo (2026).
		\bibitem{AppendixFerra} A. V. Yakushev, ``Appendix Ferra1.2: Universal Coordination Constants for Ordered and Disordered Phases,'' in \emph{YUCT}, Zenodo (2026).
		\bibitem{CoordinationSystematics} A. V. Yakushev, ``YUCT Coordination Systematics of Masses and Interactions,'' in \emph{YUCT}, Zenodo (2026).
		\bibitem{AppendixY} A. V. Yakushev, ``Appendix Y: Mathematical Foundations of YUCT and Nuclear Mass Systematics,'' in \emph{YUCT}, Zenodo (2026).
		\bibitem{AppendixOmega} A. V. Yakushev, ``Appendix $\Omega$: Cosmological Coordination and the Planck‑Scale Emergence,'' in \emph{YUCT}, Zenodo (2026).
		\bibitem{AppendixP} A. V. Yakushev, ``Appendix P: Temperature Dependence of Gravity,'' in \emph{YUCT}, Zenodo (2026).
		\bibitem{AppendixBCLM} A. V. Yakushev, ``Appendix BCLM: Biological Coordination Lagrangian,'' in \emph{YUCT}, Zenodo (2026).
		\bibitem{AppendixBCLM-EC} A. V. Yakushev, ``Appendix BCLM-EC: BCLM Coordination Clocks,'' in \emph{YUCT}, Zenodo (2026).
		\bibitem{Prigogine1977} I.~Prigogine, G.~Nicolis, \emph{Self‑Organization in Nonequilibrium Systems}. Wiley (1977).
		\bibitem{Kleiber1932} M.~Kleiber, ``Body size and metabolism,'' \emph{Hilgardia} 6, 315–353 (1932).
		\bibitem{West1997} G.~B.~West, J.~H.~Brown, B.~J.~Enquist, ``A general model for the origin of allometric scaling laws in biology,'' \emph{Science} 276, 122–126 (1997).
	\end{thebibliography}
	
\end{document}